% Topic 1.1.08 — Numerical and Functional Series.

\section{Numerical and Functional Series}
\label{sec:1.1.08-series}

\begin{ticket}
Series. Numerical and functional series. Convergence tests (d'Alembert, Cauchy, integral, Leibniz). Absolutely and conditionally convergent series.
\end{ticket}

\subsection*{Theory}

We work over the real numbers throughout this section; statements
about absolute convergence carry over verbatim to complex series by
$|z| = \sqrt{z\bar{z}}$.

\begin{definition}
\label{def:series-convergence}
A \emph{numerical series} is a formal sum
$\sum_{n=1}^{\infty} a_n$ associated with a sequence $(a_n)$ of real
numbers. Its \emph{partial sums} are
$S_N = \sum_{n=1}^{N} a_n$. The series \emph{converges} (to the
\emph{sum}~$S$) if $\lim_{N \to \infty} S_N = S$, and \emph{diverges}
otherwise.
\end{definition}

\begin{proposition}[Necessary condition]
\label{prop:necessary}
If $\sum a_n$ converges, then $a_n \to 0$ as $n \to \infty$.
\end{proposition}
\proofof{necessary}

\begin{theorem}[Cauchy criterion]
\label{thm:series-cauchy}
$\sum a_n$ converges if and only if for every $\varepsilon > 0$ there
exists~$N$ such that
$\bigl|\sum_{n=p+1}^{q} a_n\bigr| < \varepsilon$ for all $q > p \ge N$.
\end{theorem}

This is just the Cauchy criterion for the sequence $(S_N)$ in the
complete space~$\R$; we will use it without further proof.

\begin{definition}
\label{def:abs-cond-convergence}
$\sum a_n$ \emph{converges absolutely} if $\sum |a_n|$ converges.
A convergent series that does not converge absolutely is said to
\emph{converge conditionally}.
\end{definition}

\begin{theorem}[Absolute implies convergent]
\label{thm:abs-implies-conv}
If $\sum |a_n|$ converges, then $\sum a_n$ converges, and
$\bigl|\sum a_n\bigr| \le \sum |a_n|$.
\end{theorem}
\proofof{abs-implies-conv}

\medskip

\noindent\textbf{Convergence tests for series of nonnegative terms.}

\begin{theorem}[d'Alembert's ratio test]
\label{thm:dalembert}
Let $a_n \neq 0$ for all~$n$ and assume the limit
$q = \lim_{n \to \infty} \bigl|\tfrac{a_{n+1}}{a_n}\bigr|$ exists in
$[0, +\infty]$. Then
\begin{enumerate}[label=(\alph*)]
  \item if $q < 1$, $\sum a_n$ converges absolutely;
  \item if $q > 1$ (including $q = +\infty$), $\sum a_n$ diverges
        (in fact $a_n \not\to 0$);
  \item if $q = 1$, the test is inconclusive.
\end{enumerate}
\end{theorem}
\proofof{dalembert}

\begin{theorem}[Cauchy's root test]
\label{thm:cauchy-root}
Let $q = \limsup_{n \to \infty} \sqrt[n]{|a_n|} \in [0, +\infty]$.
\begin{enumerate}[label=(\alph*)]
  \item if $q < 1$, $\sum a_n$ converges absolutely;
  \item if $q > 1$, $\sum a_n$ diverges;
  \item if $q = 1$, the test is inconclusive.
\end{enumerate}
\end{theorem}
\proofof{cauchy-root}

\begin{theorem}[Cauchy--Maclaurin integral test]
\label{thm:integral-test}
Let $f \colon [1, +\infty) \to [0, +\infty)$ be a nonincreasing
nonnegative function and set $a_n = f(n)$. The series
$\sum_{n=1}^{\infty} a_n$ and the improper integral
$\int_1^{+\infty} f(x)\, dx$ converge or diverge together. Moreover,
\[
  \int_{1}^{N+1} f(x)\, dx \;\le\; \sum_{n=1}^{N} a_n \;\le\; a_1 + \int_{1}^{N} f(x)\, dx.
\]
\end{theorem}
\proofof{integral-test}

\medskip

\noindent\textbf{Alternating series.}

\begin{theorem}[Leibniz's test]
\label{thm:leibniz}
Let $(a_n)$ be a nonincreasing sequence of nonnegative numbers with
$a_n \to 0$. Then the alternating series
$\sum_{n=1}^{\infty} (-1)^{n+1} a_n$ converges, and the truncation
error satisfies
\[
  |S - S_N| \;\le\; a_{N+1}.
\]
\end{theorem}
\proofof{leibniz}

\medskip

\noindent\textbf{Functional series and uniform convergence.}

\begin{definition}
\label{def:functional-series}
Let $D \subseteq \R$ and $f_n \colon D \to \R$ for $n = 1, 2, \dots$.
The \emph{functional series} $\sum f_n$ converges \emph{pointwise on
$D$} if $\sum f_n(x)$ converges for every $x \in D$, and converges
\emph{uniformly on $D$} if
$\sup_{x \in D} \bigl|S(x) - S_N(x)\bigr| \to 0$ as $N \to \infty$,
where $S_N(x) = \sum_{n=1}^{N} f_n(x)$ and $S(x)$ is the pointwise
sum.
\end{definition}

\begin{theorem}[Weierstrass $M$-test]
\label{thm:weierstrass-m}
If there exist constants $M_n \ge 0$ such that $|f_n(x)| \le M_n$ for
every $x \in D$ and every~$n$, and if $\sum M_n$ converges, then
$\sum f_n$ converges uniformly (and absolutely) on~$D$.
\end{theorem}
\proofof{weierstrass-m}

\begin{remark}
Uniform convergence is the right notion to swap a series with limits,
derivatives, or integrals (continuity of the sum, term-by-term
integration, etc.). These applications appear in
\cref{sec:1.1.12-fourier-series,sec:1.1.13-fourier-transform}. For
deeper treatment of all results in this section see \cite{Zorich1981}.
\end{remark}

\subsection*{Proofs}

\begin{proof}[\Cref{prop:necessary}]
\proofanchor{necessary}
If $S_N \to S$, then $a_N = S_N - S_{N-1} \to S - S = 0$.
\end{proof}

\begin{proof}[\Cref{thm:abs-implies-conv}]
\proofanchor{abs-implies-conv}
Apply the Cauchy criterion (\cref{thm:series-cauchy}) to both series.
For $q > p$, the triangle inequality gives
\[
  \Bigl| \sum_{n=p+1}^{q} a_n \Bigr| \;\le\; \sum_{n=p+1}^{q} |a_n|.
\]
If $\sum |a_n|$ converges, the right-hand side is below~$\varepsilon$
for $p, q \ge N$, so the left is too, and $\sum a_n$ converges.
Passing to the limit $q \to \infty$ in
$|\sum_{n=1}^{q} a_n| \le \sum_{n=1}^{q} |a_n|$ yields the
inequality on the sums.
\end{proof}

\begin{proof}[\Cref{thm:dalembert}]
\proofanchor{dalembert}
\emph{(a) $q < 1$.} Pick $q'$ with $q < q' < 1$. By definition of the
limit, there exists~$N$ with
$\bigl|\tfrac{a_{n+1}}{a_n}\bigr| \le q'$ for all $n \ge N$. By
induction $|a_{N+k}| \le q'^k |a_N|$, so $\sum_{n \ge N} |a_n|$ is
dominated by the convergent geometric series
$|a_N| \sum_k q'^k = |a_N|/(1 - q')$. Hence $\sum |a_n|$ converges.

\emph{(b) $q > 1$.} Pick $q'$ with $1 < q' < q$. There is~$N$ with
$|a_{n+1}/a_n| \ge q'$ for $n \ge N$, so $|a_n| \ge q'^{n - N}|a_N|$
for $n \ge N$. Since $q' > 1$, this lower bound diverges, so
$a_n \not\to 0$ and the series diverges by
\cref{prop:necessary}.

\emph{(c) $q = 1$.} The harmonic series $\sum 1/n$ has $q = 1$ and
diverges, while $\sum 1/n^2$ has $q = 1$ and converges (see the
examples), so no general conclusion is possible.
\end{proof}

\begin{proof}[\Cref{thm:cauchy-root}]
\proofanchor{cauchy-root}
\emph{(a) $q < 1$.} Pick $q'$ with $q < q' < 1$. By the definition
of the limit superior, there exists~$N$ with $\sqrt[n]{|a_n|} \le q'$
for all $n \ge N$, i.e.\ $|a_n| \le q'^n$. The geometric series
$\sum q'^n$ converges, so by comparison $\sum |a_n|$ converges.

\emph{(b) $q > 1$.} There are infinitely many~$n$ with
$\sqrt[n]{|a_n|} > 1$, i.e.\ $|a_n| > 1$. Hence $a_n \not\to 0$, and
the series diverges by \cref{prop:necessary}.

\emph{(c) $q = 1$.} Same examples as for d'Alembert show
inconclusiveness.
\end{proof}

\begin{proof}[\Cref{thm:integral-test}]
\proofanchor{integral-test}
Since $f$ is nonincreasing, for each integer $n \ge 1$ and
$x \in [n, n+1]$ we have $f(n+1) \le f(x) \le f(n)$. Integrating,
\[
  a_{n+1} = f(n+1) \;\le\; \int_n^{n+1} f(x)\, dx \;\le\; f(n) = a_n.
\]
Summing for $n = 1, \dots, N - 1$:
\[
  \sum_{n=2}^{N} a_n \;\le\; \int_1^{N} f(x)\, dx \;\le\; \sum_{n=1}^{N-1} a_n.
\]
Adding $a_1$ to both sides of the left inequality gives
$\sum_{n=1}^{N} a_n \le a_1 + \int_1^N f$ — the right half of the
statement. Shifting $N \to N + 1$ in the right inequality gives
$\int_1^{N+1} f \le \sum_{n=1}^{N} a_n$ — the left half. Both
partial sums and the partial integrals are nondecreasing in~$N$ and
bound each other up to a constant; they therefore converge or diverge
together.
\end{proof}

\begin{proof}[\Cref{thm:leibniz}]
\proofanchor{leibniz}
Write $S_N = \sum_{n=1}^{N} (-1)^{n+1} a_n$. The even-indexed
partial sums satisfy
\[
  S_{2k+2} - S_{2k} = a_{2k+1} - a_{2k+2} \ge 0,
\]
so $(S_{2k})$ is nondecreasing. The odd-indexed sums satisfy
\[
  S_{2k+1} - S_{2k-1} = -a_{2k} + a_{2k+1} \le 0,
\]
so $(S_{2k+1})$ is nonincreasing. Since $S_{2k} \le S_{2k+1}$ (because
$S_{2k+1} - S_{2k} = a_{2k+1} \ge 0$), each $S_{2k}$ is bounded above
by~$S_1 = a_1$ and each $S_{2k+1}$ is bounded below by $S_2 = a_1 - a_2 \ge 0$.
Two monotone bounded sequences converge; call their limits
$L$ and $L'$. From $S_{2k+1} - S_{2k} = a_{2k+1} \to 0$, $L = L'$, and
$S_N \to L$.

For the truncation bound, both $L \in [S_{N}, S_{N+1}]$ when $N$ is
even and $L \in [S_{N+1}, S_N]$ when $N$ is odd; in either case
$|S - S_N| \le |S_{N+1} - S_N| = a_{N+1}$.
\end{proof}

\begin{proof}[\Cref{thm:weierstrass-m}]
\proofanchor{weierstrass-m}
Apply the uniform Cauchy criterion: $\sum f_n$ converges uniformly
iff for every $\varepsilon > 0$ there exists~$N$ such that
$\sup_{x \in D} \bigl|\sum_{n=p+1}^{q} f_n(x)\bigr| < \varepsilon$
for all $q > p \ge N$. Bound the supremum by the corresponding
$M$-tail:
\[
  \Bigl| \sum_{n=p+1}^{q} f_n(x) \Bigr|
  \;\le\; \sum_{n=p+1}^{q} |f_n(x)|
  \;\le\; \sum_{n=p+1}^{q} M_n.
\]
Convergence of $\sum M_n$ makes the right-hand side smaller than
$\varepsilon$ uniformly in~$x$ for $p, q$ large.
\end{proof}

\subsection*{Examples}

\begin{example}[Geometric series]
\label{ex:geometric}
For $q \in \R$, the geometric series $\sum_{n=0}^{\infty} q^n$
converges iff $|q| < 1$, in which case the sum is $1/(1 - q)$. Indeed
$S_N = (1 - q^{N+1})/(1 - q) \to 1/(1 - q)$ when $|q| < 1$, and
$|q|^{N+1} \not\to 0$ when $|q| \ge 1$. The series converges
absolutely throughout $|q| < 1$.
\end{example}

\begin{example}[Harmonic series and the $p$-series]
\label{ex:p-series}
For $p > 0$ apply the integral test to $f(x) = 1/x^p$:
$\int_1^{\infty} x^{-p}\, dx$ converges iff $p > 1$. Hence
$\sum_{n=1}^{\infty} 1/n^p$ converges iff $p > 1$. In particular the
harmonic series $\sum 1/n$ ($p = 1$) diverges, and $\sum 1/n^2$
converges.

The d'Alembert and Cauchy root tests are both inconclusive on
$p$-series ($\lim |a_{n+1}/a_n| = 1$ and
$\lim \sqrt[n]{1/n^p} = 1$), illustrating the case $q = 1$.
\end{example}

\begin{example}[$\sum 1/n!$ via d'Alembert]
\label{ex:factorial-series}
For $a_n = 1/n!$,
$\bigl|a_{n+1}/a_n\bigr| = 1/(n+1) \to 0 < 1$, so by
\cref{thm:dalembert}\,(a) the series converges absolutely. (In
fact the sum is $e - 1$, but only convergence concerns us here.)
\end{example}

\begin{example}[Conditional convergence]
\label{ex:alt-harmonic}
The alternating harmonic series
$\sum_{n=1}^{\infty} (-1)^{n+1}/n$ has $a_n = 1/n \downarrow 0$
monotonically, so it converges by \cref{thm:leibniz} (its sum is
$\ln 2$). However $\sum |a_n| = \sum 1/n$ diverges. The series
therefore converges \emph{conditionally}, not absolutely.
\end{example}

\begin{problem}
\label{prob:series-classify}
For each $\alpha > 0$ classify the series
\[
  \sum_{n=1}^{\infty} \frac{(-1)^n}{n^{\alpha}}
\]
as absolutely convergent, conditionally convergent, or divergent.
\end{problem}

\begin{proof}[Solution]
The absolute series is the $p$-series $\sum 1/n^{\alpha}$, which
by \cref{ex:p-series} converges iff $\alpha > 1$.

\emph{$\alpha > 1$.} Absolute convergence holds, hence convergence
(\cref{thm:abs-implies-conv}). The series converges absolutely.

\emph{$0 < \alpha \le 1$.} Absolute convergence fails. Apply Leibniz
to $a_n = 1/n^{\alpha}$: the sequence is nonnegative, monotonically
decreasing (since $\alpha > 0$), and tends to~$0$. Hence the
alternating series $\sum (-1)^n/n^{\alpha}$ converges by
\cref{thm:leibniz}. As absolute convergence fails, the convergence
is \emph{conditional}.

Summary: absolutely convergent iff $\alpha > 1$, conditionally
convergent iff $0 < \alpha \le 1$, never divergent for $\alpha > 0$.
\end{proof}
